\documentclass[12pt]{article}
\usepackage[utf8]{inputenc}
\usepackage[T1]{fontenc}
\usepackage[margin=1in]{geometry}
\usepackage{setspace}
\usepackage{xcolor}
\usepackage[colorlinks=true, urlcolor=blue!60!black, linkcolor=blue!60!black]{hyperref}
\setlength{\parindent}{0em}
\setlength{\parskip}{0.45em}
\renewcommand{\baselinestretch}{1.0}
\begin{document}
\begin{center}
{\Large \textit{Flast v. Cohen}, 392 U.S. 83 (1968)\par}
\vspace{0.3em}
{\normalsize Case Brief --- One-Page Edition\par}
\end{center}

\textbf{Memory Jogger.} Taxpayers may sue over congressional spending that aids religion: the double-nexus test (\href{https://supreme.justia.com/cases/federal/us/392/83/}{Justia}).

\textbf{Facts.} Florence Flast and other federal taxpayers challenged Elementary and Secondary Education Act grants that allegedly funded instruction and textbooks in religious schools, claiming standing only as taxpayers.

\textbf{Procedural History.} A three-judge court dismissed under \textit{Frothingham v.\ Mellon}, over Judge Frankel's dissent.

\textbf{Votes.} Reversed, 8--1: Warren for the Court; Douglas, Stewart, and Fortas concurring; Harlan dissenting.

\textbf{Holding.} Taxpayers have standing if they challenge an exercise of the Article I, \S\ 8 taxing and spending power as exceeding a specific limit on it, such as the Establishment Clause.

\textbf{Reasoning.} Standing turns on the party's personal stake, not separation of powers; Madison's \textit{Memorial and Remonstrance} shows the Establishment Clause limits spending. Fortas warned against a ``launching pad'' for other attacks; Harlan called taxpayers ``non-Hohfeldian'' plaintiffs.

\textbf{Examples.} \emph{Same-Side:} a statute funding church repairs. \emph{Opposite-Side:} a surplus-property gift to a seminary. \emph{Fence-Sitter:} worship funded from agency discretion within a grant program.

\textbf{Critique.} The line between congressional and executive spending, land transfers, and tax credits is arbitrary.

\textbf{Key Quotation.} ``First, the taxpayer must establish a logical link between that status and the type of legislative enactment attacked. Thus, a taxpayer will be a proper party to allege the unconstitutionality only of exercises of congressional power under the taxing and spending clause of Art. I, \S\ 8, of the Constitution.'' --- Chief Justice Warren, 392 U.S. at 102.

\end{document}
